\section{实验设置}
\subsection{数据与运行条件}

采用助教提供的 DTU scan9 测试数据，共 49 张 $1600\times1200$ 的图像，拍摄对象为一组建筑模型。数据包包含相机内外参和视图配对文件。我们依次将 49 张图像作为参考图，按配对文件的顺序选择源图，完成深度预测和点云重建。

实验在 Apple M5、24 GB 统一内存的 macOS 设备上运行，使用独立的 \texttt{mvsnet} 环境，主要软件为 Python 3.12、TensorFlow 2.16.2、NumPy 1.26.4 和 OpenCV 4.10.0。课件程序使用较早的 TensorFlow 和 CUDA，因此我们将推理接口适配到 TensorFlow 2，在 CPU 上运行；网络结构、预处理和权重均沿用上游实现，不进行训练或微调。

模型使用助教提供的 BlendedMVS \texttt{3DCNNs} 权重，代码以 MVSNet 官方仓库为基础。兼容调整后的输出已与原始推理程序核对，两者结果接近。我们保留原始数据和权重，将预测、直接生成的点云、后处理点云分别保存，方便在相同设置下比较。

\subsection{模型配置与对比方案}

输入图像按课件设置缩放至 $1152\times864$，深度图和置信度图为 $288\times216$，相机内参同步缩放。使用 192 个等间隔深度假设，间隔比例为 1.06；本数据的起始深度为 425 mm，原始间隔为 2.5 mm，因此实际采样间隔为 2.65 mm，末端深度为 931.15 mm。不启用额外的深度细化网络。

\ReportTableCaption{实验分组与参数设置}
\begin{ReportTable}{@{}p{108bp}LL@{}}
\rowcolor{ReportNavy}
\ReportHead{对比内容} & \ReportHead{变化参数} & \ReportHead{其余设置} \\
深度预测 & 输入视图数：3、5 & 分别为 1 张参考图加 2、4 张源图；权重、尺寸与深度范围相同 \\
点云生成 & 置信度阈值：0.6、0.8、0.95 & 两组预测均运行三个阈值；重点分析 5 视图组 \\
后处理对比 & 直接反投影、过滤与融合 & 使用同一组深度、相机参数和图像颜色 \\
\end{ReportTable}

点云后处理采用 CPU 上的多视角重投影检查：对每张参考图检查配对文件中的 10 张源图，要求往返像素偏差小于 1 像素、相对深度差小于 1\%，并至少有 2 张源图通过检查，即包括参考图共至少 3 个视图一致。将通过检查的深度变换回参考相机后取平均，再变换到世界坐标系，用 0.5 mm 体素合并重复点。本实现未调用依赖 CUDA 的 fusibile，因此不将两者视作完全相同的融合程序。

\subsection{统计与展示方式}

数据包未提供真实深度或参考点云，因此不计算绝对深度误差。我们比较深度图、置信度、过滤后的保留比例及点云形状。保留比例以全部参考图的像素数为分母；融合点数另计体素合并后的数量。\textbf{置信度较高、保留点较多，都不能单独证明重建更准确。}对比图统一深度色阶，点云图统一观察方向与显示比例。置信度沿用上游的四项概率和；本次有 1 个像素因重复计数超过 1，统计保留原值，显示时按 1 着色。
